\subsection{The tests in Chapter IX and a correction to Lead 3}

\begin{result}{D.1}{Proved}{The typed forms of the three tests}
Let $(M,W)$ be a relative category, let $\mathcal C=L(M,W)$, and let
$\mathscr A$ be a specified class of categories small relative to the
universe in which $M$ is being used. Write
\[
 H_{a,b}(A)=\operatorname{Hom}_{M^A[W_A^{-1}]}(a_A,b_A),\qquad
 R_A:M^A[W_A^{-1}]\longrightarrow
       \operatorname{Ho}\operatorname{Fun}(NA,\mathcal C).
\]
Here diagrams are covariant; replace $A$ by $A^{\mathrm{op}}$ for the
manuscript's coefficient convention. Rectifiability at $A$ means that
this specified functor $R_A$ is an equivalence.

At IX ms.~22--23, test (a) asks for representability of the pointed
group-valued functor denoted $\pi_1(a,b)$ in one variable, or the other,
or both: $\operatorname{Hom}(\Sigma a,b)\cong\pi_1(a,b)$ and
$\operatorname{Hom}(a,\Omega b)\cong\pi_1(a,b)$ in $M[W^{-1}]$.
These three versions must be distinguished. The second formulation of
test (b), on the bounded class $\mathscr A$, asks that
\[
 u^*:H_{a,b}(B)\longrightarrow H_{a,b}(A)
\]
be a bijection for every $u:A\to B$ in $W_\infty$ with $A,B\in\mathscr A$,
and every $a,b$. The first formulation says that the sequential path
categories are ``homotopes \`a des cat\'egories de Thomason-Kan''.
If ``homotopes'' means isomorphic in $\operatorname{Hot}$ and
``Thomason-Kan'' means Thomason fibrant, this is true of every small
category. It therefore is not a criterion. An isomorphism in
$\operatorname{Cat\,hot}$ to an object local for the indicated tests is
a different assertion; fibrancy alone cannot replace that localness.

Test (c) asks for the two natural pointed equivalences
\[
 \mathcal H(a,\Omega b)\simeq\Omega\mathcal H(a,b),\qquad
 \mathcal H(\Sigma a,b)\simeq\Omega\mathcal H(a,b),                 \tag{D.1.1}
\]
separately when the corresponding representing object exists, with
compatibility with the loop group object in $\operatorname{Hot}_\bullet$.
Here $\mathcal H$ must be a specified pointed mapping object; using the
manuscript's path object retains the comparison hypotheses of Parts B--C.
\end{result}
\begin{proof}
These are the displayed formulas and their quantifiers at IX ms.~22--23.
The phrase ``voire m\^eme des deux \`a la fois'' prevents us from silently
replacing (a) by simultaneous existence. The source itself separates the
two formulas in (c). The localization functor $NM\to\mathcal C$ sends
objectwise $W$-maps of diagrams to equivalences, so its composite with
the homotopy-category functor uniquely factors through $R_A$; this
defines its source, target, and variance. Finally, a fibrant replacement
$K\to RK$ in the Thomason model structure is in $W_\infty$ and has
fibrant target. This proves the assertion concerning the weak reading of
the first (b), without asserting a natural-transformation inverse to
$K\to RK$. The bounded formulation avoids using a single small
category as an unrestricted representative for tests of every size.
\end{proof}

\begin{result}{D.2}{Refuted}{The finite example does not fail rectifiability}
The claim in Lead 3 that the finite example at IX ms.~23 fails
rectifiability is false. It is rectifiable at every small category,
satisfies (a), both meaningful discrete versions of (b), and (c), and
its naive family satisfies Kaledin's enhanced-category axioms.
\end{result}
\begin{proof}
Let $E$ have objects $0,x$ and arrows
\[
 1_0,\quad u:0\to x,\quad v:x\to0,\quad1_x,\quad z=uv:x\to x,
 \qquad vu=1_0,\quad z^2=z\ne1_x.
\]
All remaining composites are determined by identities and these
equations; equivalently this is the full category of pointed sets on
the one- and two-element pointed sets. Set $W=\{1_0,1_x\}$.
It contains all isomorphisms, is saturated, and satisfies two out of
three. The program \texttt{d\_finite\_checks.py} checks all 34 composable
triples, the entire composition table, and the zero-object property.

For ordinary categories $A,E$, in simplicial degree $n$ there is an
equality
\[
 \operatorname{Fun}(NA,NE)_n
 =\operatorname{Hom}_{\mathrm{sSet}}(N(A\times[n]),NE)
 =\operatorname{Fun}(A\times[n],E)
 =N\operatorname{Fun}(A,E)_n.
\]
These equalities respect faces and degeneracies. Localization at
identities changes neither $E$ nor $E^A$. Consequently $R_A$ is the
canonical isomorphism for every $A$, including the empty category and
categories with nonidentity endomorphisms. This proves rectifiability,
not merely agreement on components.

For constant diagrams,
$H_{a,b}(A)=\operatorname{Hom}_E(a,b)^{\pi_0 A}$: naturality forces the
chosen arrow to be constant along every arrow, hence along each zigzag
component. This proves the second (b). The identity-backward path
calculation of Lead 2 gives a categorical homotopy equivalence with the
discrete Hom set, so the first (b) also holds with a discrete, hence local
and Thomason-fibrant, target. All based $\pi_1$ groups are trivial.
The assignments $\Sigma a=0=\Omega a$ represent them, and both sides
of (D.1.1) are contractible. The surviving distinct arrows $1_x,z$
show that the localization is not equivalent to a point.

Finally the naive family is the represented family
$A\mapsto\operatorname{Fun}(A,E)$. Kaledin's Example 7.2.1.16 applies
to every ordinary category $E$, with no finite-limit hypothesis, and
makes its contravariant version an enhanced category. Thus there is
no failed Kaledin axiom to name in this example. What it lacks is
finite products, not rectifiability or enhancement.
\end{proof}

\begin{result}{D.3}{Proved}{The actual finite-product obstruction}
A category with finitely many arrows and all finite products has at
most one arrow between any two objects. If it has a zero object as
well, it is equivalent to the terminal category.
\end{result}
\begin{proof}
Two distinct arrows $a\rightrightarrows b$ would give $2^n$ distinct
arrows $a\to b^n$ for every $n$, contradicting the fixed finite bound
on all arrows. In a pointed category a zero map exists between every
pair, so every Hom set is a singleton and every two objects are
isomorphic. The proof uses products in the category in question;
one cannot infer them from representability of the trivial
$\pi_1$-functor in D.2.
\end{proof}

\begin{result}{D.4}{Proved}{Rectifiability implies the second test (b)}
If $R_A$ is an equivalence for every $A\in\mathscr A$, then there are
natural bijections
\[
 H_{a,b}(A)\cong[BA,\operatorname{Map}_{\mathcal C}(a,b)].          \tag{D.4.1}
\]
In particular test (b) holds on $\mathscr A$. This implication uses
no two-out-of-three axiom beyond the existence of the relative
localization.
\end{result}
\begin{proof}
The space of natural transformations between the two constant
$A$-diagrams is
\[
 \operatorname{Map}_{\operatorname{Fun}(NA,\mathcal C)}
       (a_A,b_A)
 \simeq\operatorname{Fun}(NA,\operatorname{Map}_{\mathcal C}(a,b))
 \simeq\operatorname{Map}(BA,\operatorname{Map}_{\mathcal C}(a,b)).
\]
For the first equivalence, express the mapping space as the fibre of
the arrow $\infty$-category over the ordered endpoint pair.
The functor $\operatorname{Fun}(NA,-)$ preserves this limit, and the
resulting fibre is a space. For the second, a functor to a space factors
through groupoid completion, whose underlying space is $BA$.
Taking components and using the fully faithful functor $R_A$ gives
(D.4.1). The construction is natural under restriction in $A$ and in
both endpoints. Precomposition by a weak equivalence of spaces is an
equivalence of mapping spaces, hence induces the required bijection.
\end{proof}

\begin{result}{D.5}{Refuted}{The second test (b) can fail between finite shapes}
There are a six-object finite poset $X$, with all arrows weak, and a
weak equivalence $u:Y\to B$ between finite categories, such that
$u^*:H_{s,s}(B)\to H_{s,s}(Y)$ is not surjective.
Thus the formulation of (b) using only weak equivalence to a
Thomason-fibrant category is not equivalent to its constant-diagram
formulation, even with finite test categories.
\end{result}
\begin{proof}
Let $Y$ have minima $r_0,r_1$, maxima $q_0,q_1$, and all four
relations $r_i<q_j$. Adjoin to $Y$ two incomparable objects $s,n$,
both above every object of $Y$, to obtain $X$. These relations and
transitivity give all arrows. The nerves realize respectively a
circle and its unreduced suspension, a $2$-sphere. Let $B$ be the
category with objects $a,b$ and exactly two nonidentity arrows
$f,g:a\rightrightarrows b$; there are no composable pairs of
nonidentity arrows. Send $r_i$ to $a$, $q_j$ to $b$, the edge
$r_0<q_1$ to $g$, and the other three edges to $f$.
The oriented cycle
\[
 \gamma=(r_0q_0)-(r_1q_0)+(r_1q_1)-(r_0q_1)
\]
maps to $f-g$. Both realizations are circles and $Bu$ has degree one,
so $u\in W_\infty$.

A functor $B\to X$ must send $f$ and $g$ to the same arrow because
$X$ is a poset. This identifies the entire functor category, including
natural transformations, with $X^{[1]}$. Evaluation at the source is
a homotopy equivalence $BX^{[1]}\to BX$: the constant-arrow section
has a natural transformation to the identity. Since $BX\simeq S^2$,
\[
 H_{s,s}(B)=\operatorname{Aut}_{\Pi_1 B(X^{[1]})}(s_s)=1.
\]
Here localization at all arrows is the fundamental groupoid of the
nerve; this statement concerns ordinary localization, not an assumed
diagram-rectification theorem.

Write $j:Y\hookrightarrow X$ and choose $y_0=r_0$. In $X^Y$ the
natural-transformation loop
\[
 s_Y\longleftarrow j\longrightarrow n_Y
       \longleftarrow(y_0)_Y\longrightarrow s_Y                 \tag{D.5.1}
\]
evaluates, on $BY\times S^1$, to the fundamental suspension cycle
$n*\gamma-s*\gamma$. The two varying meridians give its two cones;
the fixed return meridian contributes no two-simplex. This is a
nonzero integral $2$-cycle, and $NX$ has no nondegenerate
$3$-simplices. Its coefficient on each of the eight triangles is
$+1$ or $-1$, with zero boundary. Evaluation and pairing with the
fundamental cohomology class of $BX$ therefore define a homomorphism
from $\pi_1 B(X^Y)$ at $s_Y$ to $\mathbb Z$ which sends (D.5.1) to
$1$. Its powers are distinct. This proves that $H_{s,s}(Y)$ contains
an infinite-order element, proving nonsurjectivity.

The attached program checks the posets and relations, the image of
$\gamma$, all eight coefficients and their zero boundary, and the
absence of three-simplices. It finds 198 objects of $X^Y$ and 18
objects of $X^B$. This sharpens the companion's ms.~134 calculation
by replacing its infinite category $B\mathbb Z$ by the finite
parallel-arrow category $B$.
\end{proof}

\begin{result}{D.6}{Proved}{What follows formally from the higher loop test}
Each corresponding formula in (c) implies the associated
representability in (a). In a pointed $\infty$-category with the
required homotopy pullbacks and pushouts, the choices
$\Omega b=0\times_b0$ and $\Sigma a=0\amalg_a0$ satisfy (c), with
all loop coherences. Test (b) on finite posets alone does not force
the existence of both representing functors.
\end{result}
\begin{proof}
Apply $\pi_0$ to (D.1.1), using
$\pi_0\Omega K=\pi_1(K,*)$. Conversely, the universal property of
the displayed pullback gives
$\operatorname{Map}(a,0\times_b0)\simeq
 *\times_{\operatorname{Map}(a,b)}*$, and the pushout gives the
dual formula. These equivalences use the same iterated pullbacks
that define loop multiplication, so they respect its coherent
structure; the two representabilities yield the adjunction by Yoneda.

For the last assertion use finite pointed simplicial sets with weak
homotopy equivalences. This is a finite cofibration category, as in
D.11 below, and hence is rectifiable at finite posets by D.7. Thus
(b) holds there by D.4. A representing object for $\Omega S^1$
cannot be finite: evaluation at the pointed $0$-sphere would give
\[
 \operatorname{Hom}(S^0,\Omega S^1)=
 \pi_1\operatorname{Map}_*(S^0,S^1)=\mathbb Z,
\]
whereas the left side is the finite set $\pi_0$ of any finite
candidate. The opposite relative category supplies the dual
failure for suspension. The equivalence of its localization with
the full subcategory of finite CW homotopy types follows directly
from simplicial approximation: maps and homotopies between finite
polyhedra are represented after finite subdivision, which is a weak
equivalence, and subdivision of any finite homotopy is still finite.
No conclusion that the disjunctive version of (a) fails is being
drawn from this example.
\end{proof}

\subsection{Sufficient hypotheses and intrinsic cofibrations}

\begin{result}{D.7}{Known}{Rectification theorems with their actual shape bounds}
The following sufficient conditions concern the canonical comparison
$R_A$, and in fact give equivalences before taking homotopy categories.
An ordinary category with weak equivalences and fibrations in the
sense of Cisinski rectifies every finite inverse category; its dual
rectifies every finite direct category. In particular each rectifies
every finite poset. An ABC cofibration category satisfying all six
axioms, or its dual, rectifies all categories small relative to the
universe of its coproduct and sequential-colimit axioms. The same is
true for a Quillen model category complete and cocomplete relative
to that universe.
\end{result}
\begin{proof}
For finite diagrams the exact localization statement is
\citet[Theorem 7.6.17]{DCisinski2019}:
\[
 L(\operatorname{Fun}(N I^{\mathrm{op}},C))
       \longrightarrow\operatorname{Fun}(N I^{\mathrm{op}},L C)
\]
is an equivalence for finite direct $I$ and an $\infty$-category with
weak equivalences and fibrations. Apply it to $C=NM$; nerve preserves
ordinary diagram categories. Dualizing gives the covariant
cofibration assertion. Definition 7.4.12 of that source allows
nonfibrant objects and does not assume functorial factorizations.
The ordinary finite theory is developed in
\citet[Definition 1.1, Theorems 1.30 and 2.21, and Corollary
3.20]{DCisinski2010}; its coefficient convention is contravariant.

For infinite shapes use \citet[Theorem 3.33]{DLenz2018}, together
with the localization comparison for quasicategories of frames.
The word ``cofibration category'' there includes arbitrary small
coproducts and countable composites: Definition 1.5 and Warning 1.7
must be retained. An equivalent canonical formulation is
\citet[Theorem 7.9.8]{DCisinski2019}, with the hereditary hypothesis
in Definition 7.9.1; Remark 7.9.7 explains the sufficient
sequential-colimit axiom. These are statements about localization
of the diagram category, not merely about existence of an abstract
equivalent derivator.

For a general model category, its cofibrant objects satisfy the ABC
axioms in the stated universe. The cofibrant replacement comparisons
for diagram categories used in Lenz's proof identify the induced
functor with $R_A$. The underlying derivator theorem is also
\citet[Corollaire 6.21]{DCisinski2010},
and the model-category case is treated in
\citet[Th\'eor\`eme 6.11]{DCisinski2003}. The result is therefore not restricted to
combinatorial or simplicial model categories. It does not assert
rectification at diagrams larger than the universe for which the
model category has these limits and colimits.
\end{proof}

\begin{result}{D.8}{Proved}{Enough intrinsic cofibrations give finite rectification}
Suppose $M$ has an initial object and all pushouts, $W$ contains the
isomorphisms and satisfies two out of three, and every map factors
as an intrinsic $W$-cofibration followed by a member of $W$.
Then $(M,\operatorname{Cof}_W,W)$ satisfies R\u adulescu-Banu's
precofibration axioms CF1--CF4 \citep[Definition 1.1.1]{DRadulescu2009},
equivalently the dual of Cisinski's
derivable-category axioms. Consequently $(M,W)$ is rectifiable at
every finite direct category, and in particular at every finite
poset and every span. Thus, under these hypotheses including the
initial object, no counterexample asked for in the POLES addendum
can exist at a finite poset.
\end{result}
\begin{proof}
An intrinsic cofibration $i:A\to B$ means that the pushout functor
$i_!:A\backslash M\to B\backslash M$ preserves underlying
$W$-maps. Isomorphisms have this property. Two successive pushouts
identify with pushout along the composite, so these maps compose;
the same identification shows stability under arbitrary pushout.
The initial object's identity is a cofibration. These are CF1.
CF2 is the assumed two-out-of-three property, and CF4 is the given
factorization, stronger than its restriction to cofibrant sources.

It remains to prove that a trivial intrinsic cofibration remains
weak after arbitrary cobase change. Let $i:A\to B$ lie in both
classes and factor a proposed base-change map as
$A\xrightarrow{j}A'\xrightarrow{w}A''$ with $j$ intrinsic and
$w\in W$. Put $B'=B\amalg_A A'$ and
$B''=B'\amalg_{A'}A''$. Since $j$ is intrinsic and $i\in W$,
$A'\to B'$ is in $W$. Since it is also intrinsic,
$B'\to B''$ is in $W$. In the commuting square with the other
vertical arrow $w:A'\to A''$, two out of three now gives
$A''\to B''\in W$. This proves CF3, with actual comparison maps.
It is precisely the needed part of POLES A.9, verified here from
its definition rather than imported from a gluing theorem.

D.7 now applies to the nerve of this category with weak
equivalences and cofibrations. It includes the replacement of
noncofibrant objects, so we have not assumed that every
$0\to X$ is an intrinsic cofibration. In Szumi\l o's convention \citep[Section 1.1, axioms C0--C5]{DSzumilo2014},
one additionally restricts to cofibrant objects and requires the
stated two-out-of-six condition; the stronger convention is not
needed to invoke Cisinski's theorem. In Lenz's convention one
must also verify CF5--CF6. Neither change of convention is implicit
in the finite conclusion.
\end{proof}

\begin{result}{D.9}{Refuted}{Initial objects do not make every object intrinsically cofibrant}
The hypotheses of D.8 do not imply that every object is
$\operatorname{Cof}_W$-cofibrant. They also do not imply the infinite
coproduct and sequential-colimit axioms CF5--CF6.
\end{result}
\begin{proof}
Take the diamond lattice $0<a,b<1$ and
$W=\{\text{identities},0\to a,0\to b\}$.
All pushouts are joins and $0$ is initial. The two-out-of-three
axiom is checked on its only nontrivial composable pairs
$0<a<1$ and $0<b<1$. The complete intrinsic class consists of
the identities and $0\to1,a\to1,b\to1$. Indeed, pushout of
$0\to b\in W$ along $0\to a$ is $a\to1\notin W$, and the
dual argument excludes $0\to b$. Every map into $1$ is intrinsic
because its pushout sends every object under its source to $1$.
Every weak map factors as its identity followed by itself; all
other nonidentity maps are intrinsic and factor followed by an
identity. This verifies every hypothesis of D.8 but excludes
cofibrancy of $a,b$. The complete finite enumeration is attached.

For the infinite axioms take finite sets, with $W$ the
isomorphisms. Every map is intrinsic, every factorization can end
in an identity, and finite pushouts exist. A countable coproduct
of singletons does not exist. Neither does the colimit of the
cofibration sequence $\{0\}\subset\{0,1\}\subset\cdots$.
To verify the latter without relying on a set-theoretic picture,
a putative finite colimit $T$ would give
$\operatorname{Hom}(T,\{0,1\})\cong\{0,1\}^{\mathbb N}$,
which is impossible. Thus both CF5 and CF6 can fail. This
example remains rectifiable, since its weak equivalences are
isomorphisms; the failure concerns the proposed axioms, not the
conclusion of an unproved converse.
\end{proof}

\subsection{Counterexamples separating the proposed criteria}

\begin{result}{D.10}{Refuted}{All cubes do not detect rectification at finite posets}
Rectifiability at $[1]$ and all its finite products does not imply
rectifiability at all finite posets. A counterexample has six objects
and twelve arrows, including identities.
\end{result}
\begin{proof}
Let $P$ be the face poset of the boundary of a triangle, with objects
$v_0,v_1,v_2,e_{01},e_{12},e_{20}$, relations given by incidence,
and $W=\operatorname{Mor}(P)$. Its nerve realizes to a circle, so
$L(P,W)\simeq B\mathbb Z$, regarded as an $\infty$-groupoid.
For any category $A$ with an initial object $a_0$, evaluation
$\operatorname{Fun}(A,P)\to P$ has a section by constants and a natural
transformation from the section followed by evaluation to the identity.
Consequently it induces an equivalence after groupoid completion.
Evaluation also gives an equivalence
$\operatorname{Fun}(NA,BP)\simeq BP$, since $BA$ is contractible.
The square of these two evaluation maps with $R_A$ commutes. Thus
$R_A$ is an equivalence. This applies to every cube, including $[0]$.

For $A=P$, the objects on the strict side are monotone maps $P\to P$.
The degree of every such map is $-1,0$, or $1$: a simplicial map of the
six-edge polygon cannot traverse more than six edges, whereas a map of
degree $d$ must traverse at least $6|d|$ edges after cancellations.
The target $\operatorname{Ho}\operatorname{Fun}(NP,BP)$ has objects
of every integral degree. Isomorphic objects there have equal degrees,
so the degree-two object is not in the essential image. The exhaustive
enumeration in \texttt{d\_finite\_checks.py} finds 234 monotone maps,
with degree counts $3,228,3$ for $-1,0,1$, respectively. This verifies
the entire finite search, independently of the length argument.
\end{proof}

\begin{result}{D.11}{Proved}{Finite posets versus the infinite cyclic group}
There is a small relative category rectifiable at every finite direct
category, hence every finite poset, but not at $B\mathbb Z$.
There is such an example fibrant in the usual Barwick--Kan model
structure. Thus Barwick--Kan fibrancy does not imply rectification at
all small shapes.
\end{result}
\begin{proof}
Use a small skeleton $M$ of finite simplicial sets, where finite means
having finitely many nondegenerate simplices. Give it weak homotopy
equivalences and monomorphisms as cofibrations. Finite pushouts remain
finite; a pushout of a trivial monomorphism is a weak equivalence.
For a map $f:A\to B$, its finite simplicial mapping cylinder factors
$f$ as a monomorphism followed by a simplicial homotopy equivalence.
All objects are cofibrant. Thus the finite cofibration axioms hold,
and D.7 applies. The same statements hold for pointed finite
simplicial sets, used in D.6.

Any strict functor $B\mathbb Z\to M$ consists of a finite simplicial
set $K$ with an actual simplicial automorphism $t$. Its automorphism
group embeds in the permutation group of the finite set of
nondegenerate simplices. Therefore $t$ has finite order. In contrast,
let $X$ be a finite simplicial model of a wedge of two circles.
The Nielsen homotopy equivalence sending the free generators
$a\mapsto a$, $b\mapsto ab$ induces, in a chosen basis of $H_1$,
\[
 T=\begin{pmatrix}1&1\\0&1\end{pmatrix},\qquad
 T^n=\begin{pmatrix}1&n\\0&1\end{pmatrix}.
\]
It gives an infinite-order automorphism of $X$ in $L(M,W)$.
To see that this automorphism belongs to the stated localization,
triangulate its graph map after a finite subdivision and invert the
subdivision map; its homotopy inverse is represented the same way.
Equivalently, the usual finite cylinder homotopy calculus gives the
homotopy category of finite CW types. A loop representing this
equivalence in the core of $L(M,W)$ defines a coherent functor
$B\mathbb Z\to L(M,W)$.

An equivalence of coherent $B\mathbb Z$-diagrams conjugates their
monodromies in the homotopy category. After applying $H_1$, finite-order
monodromy remains finite order under conjugacy. It cannot be conjugate
to $T$. This proves failure of essential surjectivity of $R_{B\mathbb Z}$.
The matrix computation is included in the executable check; its
displayed general formula proves infinite order without a finite-test
inference.

Finally take $M^{\mathrm{op}}$. This is a fibration category and is
Barwick--Kan fibrant by \citet[Theorem 4.13]{DMeier2016}. Localization
and diagram comparison commute with opposites, with the indexing
category also replaced by its opposite. Since
$B\mathbb Z\cong(B\mathbb Z)^{\mathrm{op}}$, the same obstruction
applies. For $M$ itself the cited theorem gives fibrancy in the
\emph{conjugate} Barwick--Kan structure (Corollary 4.15); it is essential
to distinguish the two structures.
\end{proof}

\begin{result}{D.12}{Refuted}{Rectification does not imply Barwick--Kan fibrancy}
A nine-object relative poset is rectifiable at every small shape but
is not Barwick--Kan fibrant.
\end{result}
\begin{proof}
Let $F$ be the fence
\[
 0<1>2<3>4<5>6<7>8,
 \qquad W=\operatorname{Mor}(F).
\]
Successively delete the left endpoint of the remaining fence, sending
it to its only neighbour. At each step the retraction followed by
inclusion is comparable pointwise to the identity, hence admits a
natural transformation in one direction. The eight steps give a fixed
finite zigzag from $\operatorname{id}_F$ to a constant functor. Applying
$\operatorname{Fun}(A,-)$ gives such a contraction for every $A$, with
no dependence of the number of steps on $A$. Thus both
$L(F,W)$ and $L(\operatorname{Fun}(A,F),W_A)$ are contractible spaces.
The specified comparison $R_A$ is an equivalence.

For a relative category with every arrow weak, the zeroth level of the
Barwick--Kan bisimplicial nerve is
\[
 (N_\xi F)_{0,\bullet}=\operatorname{Ex}^{2}NF.
\]
This follows directly from double subdivision and
$\xi\widehat{[n]}\cong c\operatorname{Sd}^{2}\Delta^n$ on underlying
categories; see \citet[Example 2.5 and Definition 2.6]{DMeier2016}.
An edge of this simplicial set is a monotone map from the five-vertex
fence $0<1>2<3>4$ to $F$, with the indicated endpoint values.
There are edges $0\to4$ and $4\to8$, obtained from the consecutive
subfences. There is no edge $0\to8$: every successive pair of values
is equal or adjacent in the undirected nine-vertex fence, so four
steps cannot cover distance eight. These two edges define an inner
two-horn whose filler would supply the missing edge. Hence
$\operatorname{Ex}^{2}NF$ is not Kan. A fibrant relative category has
a Reedy fibrant $N_\xi$, so its zeroth level must be Kan. This proves
the claim. The program enumerates every five-fence map and verifies
both the missing edge and all eight contraction retractions.
\end{proof}

\begin{result}{D.13}{Refuted}{A relative poset presentation need not rectify}
Presenting an ordinary category as an $\infty$-localization of a
relative poset does not imply that the naive strict diagram
localizations are its coherent diagram categories. The presenting
relative poset can be finite.
\end{result}
\begin{proof}
Use the six-element $P$ of D.10. Choose an oriented circuit in its
nerve. Define $q:P\to B\mathbb Z$ to send one incidence edge to
the generator (or its inverse according to orientation) and all
other incidence edges to $0\in\mathbb Z$. There are no composable
nonidentity pairs, so these assignments define a functor. Its map
on classifying spaces has degree one. Since every arrow maps to an
isomorphism and $BP$ is a circle, $q$ exhibits $B\mathbb Z$ as the
$\infty$-localization of $(P,\operatorname{Mor}(P))$.

Every strict $B\mathbb Z$-diagram in a poset is constant: its
generator must act by an automorphism, and posets have only identity
automorphisms. Therefore $P^{B\mathbb Z}\cong P$. Its localization
is a connected groupoid. On the coherent side,
$\operatorname{Fun}(B\mathbb Z,B\mathbb Z)$ has components indexed
by the homomorphisms $\mathbb Z\to\mathbb Z$, that is, by integers.
The comparison lands in the component of the trivial homomorphism.
It is not essentially surjective. It already fails at the finite
poset $P$ by D.10.

There is also a presentation of a \emph{finite} category for which
rectification fails at that same finite category. Let $K$ be the
infinite join of copies of $S^0$, indexed by $\mathbb N$. Its
nonempty faces are finite sets $F\subset\mathbb N\times\{+,-\}$
containing at most one sign at each index. Simultaneous sign reversal
acts freely. Let $Q$ be the poset of its face orbits, with
$[F]\leq[G]$ if $F\subseteq G$ after a choice of signs. Nonemptiness
makes the aligning sign unique. A chain of face orbits therefore
lifts uniquely after one representative is chosen, so
\[
 NQ\cong N(\operatorname{Face}K)/C_2,
 \qquad BQ\simeq\mathbb{RP}^{\infty}\simeq BC_2.
\]
Here $K$ is contractible: its finite joins are spheres of dimensions
tending to infinity, so every homotopy group of the union vanishes;
the CW Whitehead theorem applies. The action is free, and the
quotient is the displayed classifying space.

Normalize each face representative so that its least-index sign is
$+$. For $[F]\leq[G]$ let $q([F]\leq[G])\in C_2$ be the unique
sign $\varepsilon$ with $\varepsilon F\subset G$. Signs multiply
under composition, so $q:Q\to BC_2$ is a functor. It classifies
the free cover of face representatives and induces the displayed
equivalence after groupoid completion. All arrows are marked.
But $Q^{BC_2}\cong Q$, whereas
$\operatorname{Fun}(BC_2,BC_2)$ has two components, indexed by
the trivial and identity homomorphisms of $C_2$. The canonical
comparison misses the identity component. The target category and
the test shape each have one object and two arrows; their complete
composition law and homomorphisms are included in the finite check.

The general existence of a relative poset presentation is
\citet[Theorem 6.3.6.1, tag 02MD]{ALurie2026}. Its universal property
concerns \emph{coherent} functors out of the presenting poset that
invert the marked arrows. It does not identify localization of the
ordinary category of strict functors into that poset. The two
variances and the placement of the localization are different.
\end{proof}

\subsection{Kaledin's intrinsic axioms and the localization they represent}

\begin{result}{D.14}{Known}{The precise shapes and axioms in Kaledin v3}
In \citet{DKaledin2026}, $\mathrm{Pos}$ means small posets of finite
chain dimension; it does not mean finite posets. The notation
$\mathrm{pos}$ denotes finite posets. The category $\mathrm{Pos}^{+}$
consists of left-bounded posets, and $\mathrm{Pos}^{\pm}$ of left-finite
posets. Left-bounded means that every principal lower set has finite
chain dimension; left-finite means that every such set is finite.

A restricted enhanced category is a small additive, semiexact,
separated, nondegenerate reflexive family over $\mathrm{Pos}$,
satisfying excision and the cylinder axiom. An enhanced category is
the analogous family over $\mathrm{Pos}^{+}$ with semicontinuity
added. These are Definitions 7.2.1.1, 7.2.1.3 and 7.2.1.4, using
Definitions 7.1.1.1--7.1.1.2, 7.1.1.9 and 7.1.4.6; the shapes are
Definitions 2.1.5.3--2.1.5.4. The representing complete Segal space in
Theorem 7.3.1.3 can be chosen $\kappa$-small when the family is
$\kappa$-bounded for an uncountable regular $\kappa$.
\end{result}
\begin{proof}[Source audit and interpretation]
The axioms have distinct content. Nondegeneracy says that evaluation
at all points is conservative. Reflexivity specifies the adjunctions
associated with $J\to J\times[1]\to J$. Separatedness asks that
the induced functor $C_{J\times[1]}\to\operatorname{Fun}([1],C_J)$
be full and essentially surjective. Additivity identifies the value
on a disjoint union with the product of the values, including infinite
unions allowed by the shape universe. Semiexactness requires an
epivalence for the prescribed standard pushout square. Excision,
cylinder and semicontinuity require semicartesianness for precisely
the squares (6.3.2.1), (6.3.2.3), (6.3.2.2), respectively. They are
conditions on the family and its transition functors, not conditions
on its value at the point alone.

Proposition 7.2.3.7 extends a \emph{good reflexive family} on
$\mathrm{Pos}^{\pm}$ to a bar-invariant ample subcategory containing
$\mathrm{Pos}^{\pm}$; it preserves the additional listed axioms.
It does not extend an arbitrary family on finite posets. In Lemma
7.2.4.1 the naive localization of a model category on left-finite
posets has an enhanced extension; its identification with the naive
family on $\mathrm{Pos}^{+}$ is asserted when the model category is
complete. This extra clause must not be discarded. Chapters 6 and 7
of v3, dated 23 February 2026, remain marked preliminary by the author;
we treat the stated representation theorem as a preprint theorem,
not as a separately re-proved or refereed result.
\end{proof}

\begin{result}{D.15}{Known}{Intrinsic recognition on Kaledin's shape category}
Fix compatible universes and a small relative category $(M,W)$.
Let $\mathcal D_J=M^{J^{\mathrm{op}}}[W_J^{-1}]$ with its actual
restriction functors. The following are equivalent:
\begin{enumerate}
\item This family on $\mathrm{Pos}$ satisfies the restricted
      enhanced-category axioms in D.14.
\item The canonical comparisons
\[
 \mathcal D_J\longrightarrow
 \operatorname{Ho}\operatorname{Fun}(NJ^{\mathrm{op}},L(M,W))
\]
are equivalences for every $J\in\mathrm{Pos}$.
\end{enumerate}
The analogous equivalence holds on $\mathrm{Pos}^{+}$ with the full
enhanced-category axioms. In particular, the category represented by
a naive enhanced family is the specified localization $L(M,W)$,
not an unidentified $\infty$-category. This is a formal consequence
of Kaledin's theorem and the localization universal property; it is
not claimed as a new independent recognition theorem.
\end{result}
\begin{proof}
Write $\mathcal K(X)_J=\operatorname{Ho}\operatorname{Fun}
(NJ^{\mathrm{op}},X)$, and use the associated Grothendieck fibration.
The identification of this notation with Kaledin's $K(X)$ follows
from the expansion in Lemma 7.3.1.1, formulas (7.3.1.7)--(7.3.1.8):
its core at $J$ is the fundamental groupoid of the derived space
of $NJ^{\mathrm{op}}$-diagrams, and expansion recovers their
homotopy category. For ordinary $M$, this is exactly the represented
family $K(M)$ of (7.1.1.5) and Example 7.2.1.16.

We need the following elementary two-categorical localization fact.
The pointwise localization map
\[
 \ell:K(M)\longrightarrow\mathcal D
\]
has fully faithful precomposition on categories of cartesian
functors into any category-valued family. In each fibre this is
the ordinary localization universal property, including natural
transformations: a transformation between two functors on the
localized category is uniquely determined on the original category,
and naturality for inverses follows from naturality for the inverted
arrows. The compatibility isomorphisms for restrictions descend in
the same way. Their coherence equations hold because they can be
checked before localization. This proves the asserted fact for
cartesian functors and modifications.

Assume (1). Theorem 7.3.1.3 identifies $\mathcal D$ with $K(X)$ for
a complete Segal space $X$. By Definition 7.2.2.1 an enhanced functor
is exactly a cartesian functor over the shape category; it need not
preserve Kan extensions. The same theorem, applied to $\ell$, gives
$f:NM\to X$ with $K(f)\cong\ell$. At the point, every $w\in W$
is sent to an isomorphism in $\operatorname{Ho}X$, so $f$ inverts
$W$. Put $C=L(M,W)$, with $q:NM\to C$. Its universal property gives
$g:C\to X$ and an equivalence $gq\simeq f$.

The canonical diagram comparison $\theta:\mathcal D\to K(C)$
commutes with restriction, hence is an enhanced functor. The theorem
represents it by $h:X\to C$, with $hf\simeq q$. Consequently
$hgq\simeq q$. Precomposition with $q$ is fully faithful on
$\infty$-functor categories, so $hg\simeq\operatorname{id}_C$.
On the other side,
\[
 K(g)\theta\ell\cong K(gq)\cong K(f)\cong\ell.
\]
The pointwise localization fact proved above gives
$K(g)\theta\cong\operatorname{id}_{\mathcal D}$. The last clause
of Theorem 7.3.1.3, which identifies homotopies with isomorphisms of
enhanced functors, now gives $gh\simeq\operatorname{id}_X$.
Thus $h$ is an equivalence, and its represented functor is the
\emph{specified} comparison $\theta$. This proves (2).

Conversely the canonical comparisons form a cartesian equivalence
of families. The coherent family $K(C)$ is enhanced by Lemma 7.3.1.1,
and the axioms are invariant under such equivalences. The argument
for $\mathrm{Pos}^{+}$ is identical, using the corresponding clauses
of Theorem 7.3.1.3. No two-out-of-three assumption is used in this
formal argument beyond whatever relative-category convention is
chosen for forming the localization.
\end{proof}

\begin{result}{D.16}{Proved}{An explicit infinite-poset failure of additivity}
For the six-object relative poset $(P,\operatorname{Mor}(P))$ of D.10,
the naive family fails Kaledin's additivity axiom at the countable
discrete poset $I=\mathbb N$. This shape has chain dimension zero
and is left-finite. In contrast, the finite category $E$ of D.2
satisfies every enhanced-category axiom: there is no failed axiom
to assign to that example.
\end{result}
\begin{proof}
Since every arrow is weak, the fibre at $I$ is
$\Pi_1 B(P^I)$. Choose a vertex $a\in P$ and consider the comparison
on its constant object:
\[
 \operatorname{Aut}_{\Pi_1 B(P^I)}(a^I)
       \longrightarrow
 \prod_{n\in\mathbb N}\operatorname{Aut}_{\Pi_1 BP}(a)
       =\mathbb Z^{\mathbb N}.
\]
Every morphism on the left is represented by a finite zigzag of
arrows in the product category. If it has length $r$, every
coordinate loop traverses at most $r$ edges of the six-edge circle.
Its coordinate degrees are therefore uniformly bounded in absolute
value, for example by $r$. The element $(n)_{n\in\mathbb N}$ of
the product cannot be in the image. Thus the additivity functor
$\mathcal D_I\to\prod_I\mathcal D_{\mathrm{pt}}$ is not full.
This is exactly Definition 7.1.4.6, not merely a failure of an
unspecified enhancement property. It also exhibits failure of
rectifiability at this particular infinite poset.

For $E$ and identity weak equivalences, the family is $K(E)$.
Example 7.2.1.16 states and verifies every axiom for $K(E)$ for
\emph{any} ordinary category $E$, by the fact that the prescribed
squares are pushouts in $\mathrm{Cat}$. Finite products or finite
limits in $E$ are not required. The failure of a product $x\times x$
in D.3 therefore concerns completeness, not enhancement.

The extension theorem cannot repair this example while retaining its
naive values on all left-finite posets, since $I$ already belongs to
that class. More generally, when a good left-finite family has an
enhanced extension, uniqueness of that extension does not assert
that an independently defined naive family on larger posets is the
extension. D.15 characterizes when the latter identification holds.
\end{proof}

\subsection{The later manuscript classes and the remaining boundaries}

\begin{result}{D.17}{Proved}{What the later cofibration and Quillen axioms imply}
The following implications hold, with the displayed size assumptions.
\begin{enumerate}
\item The axioms Cof$_0$--Cof$_3$ of XVII ms.~39--40 and XIX
      ms.~33--36, with the stated initial object and the pushouts
      required by the axioms, give rectification at every finite
      direct category. They do not require every object to be
      cofibrant.
\item A quill\'enisable pair in the sense of XIII ms.~50--51 is
      rectifiable at finite direct categories. Already M0, M2,
      M$'$ and M$''$ at XIII ms.~7--10 supply enough intrinsic
      cofibrations and enough intrinsic fibrations.
\item Adding the full transfinite axiom Cof$_4$ of XIX ms.~33,
      including existence of the indicated small filtered colimits,
      gives the ABC coproduct and sequential-colimit axioms.
      Rectification then holds at all shapes small in that universe,
      and the left Kan extensions asked for in XIX Question 3.14,
      ms.~48, exist. No small generating set or functorial
      factorization is needed for this conclusion.
\item More specifically, enough intrinsic cofibrations, an initial
      object, pushouts, small filtered colimits, and preservation of
      $W$ by such filtered colimits imply (3).
\end{enumerate}
\end{result}
\begin{proof}
For (1), Cof$_0$ is two out of three with isomorphisms; Cof$_1$ is
closure of cofibrations under composition and cobase change;
Cof$_2$ is closure of trivial cofibrations under cobase change;
Cof$_3$ factors every map as a cofibration followed by a weak
equivalence. Together with the initial object these are CF1--CF4,
so D.7 applies. The source at XVII ms.~40 explicitly restricts the
Brown category to the full subcategory of cofibrant objects and
compares its localization with that of the original category.
The same replacement, included in Cisinski's theorem, avoids an
unjustified all-objects-cofibrant assumption.

For (2), M0 supplies finite limits and colimits. One M2 factorization
has first map in $C\subset\operatorname{Cof}_W$ by M$''$ and second
in $TF\subset W$ by M$'$. The other has first map in
$TC\subset W$ and second in $F\subset\operatorname{Fib}_W$.
Apply D.8 and its dual. The definition of ``exact weak Quillen
structure'' further asks for $W=TF\circ TC$ (M2 bis); it must not
be substituted without proof for the stronger identities
$TC=C\cap W$, $TF=F\cap W$. For finite rectification this extra
exactness is unnecessary. Likewise the word ``quill\'enisable'' by
itself does not supply colimits of unbounded diagrams. XIII
ms.~53 explicitly discusses examples without small filtered
colimits.

For (3), the initial object and finite pushouts give finite
coproducts. A small coproduct is the filtered colimit of its finite
subcoproducts, so it exists. To verify CF5 for a family of
cofibrations, well order the family and replace one summand at each
successor step. Each step is a cobase change of that cofibration;
at a limit take the prescribed filtered colimit. Transfinite closure
of $C$ proves that the resulting coproduct map is a cofibration.
If every map was trivial, use transfinite closure of $C\cap W$.
Coproducts of cofibrant objects are cofibrant by the same argument
starting at the initial object. Cof$_4$ applied to a countable
sequence supplies CF6, including its assertion for trivial
cofibrations. Hence D.7 applies in the ABC convention, including
the associated left derivator and its left Kan extensions.

For (4), a transfinite composite of intrinsic cofibrations is
intrinsic. Given a weak map under its source, push it out at each
stage. At successor stages it remains weak by the defining property
of the relevant intrinsic cofibration; at limit stages it remains
weak by the assumed filtered-colimit stability. Pushout commutes
with the colimit, so the last comparison is precisely the required
pushout map. If each stage is also weak, the transfinite composite
is weak by the same induction, applying filtered-colimit stability
to the diagram from the constant initial stage. D.8 identifies
trivial intrinsic cofibrations and supplies the remaining finite
axioms. This proves Cof$_4$ and then (3).

These conclusions also locate two overstatements to avoid. The
Brown assertion at XIX ms.~35 about ambient intrinsic cofibrations
must be read with the restriction explicitly present at XVII
ms.~40, or with the additional properness hypothesis
$C\subset\operatorname{Cof}_W$. Further, the functorial
factorization in XIX Corollary 3.3 alone gives objectwise
factorizations of diagram categories; it does not add CF5--CF6.
Our all-shape conclusion uses their verified transfinite source,
not functoriality as a replacement for infinite colimits.
\end{proof}

\begin{result}{D.18}{Known}{Accessible localizers on categories}
For every accessible fundamental localizer $W$, the pair
$(\operatorname{Cat}_{\mathcal U},W)$ is rectifiable at all
$\mathcal U$-small categories and has the associated strong
derivator. It has enough intrinsic cofibrations. Enough intrinsic
fibrations, in contrast, is equivalent to right properness of its
Thomason model structure, or to generation of $W$ by maps $A\to e$.
The pointed versions of (a) and (c) hold in its pointed model
category; (b) holds before pointing as well.
\end{result}
\begin{proof}
\citet[Th\'eor\`eme 5.2.15]{DCisinski2006} constructs the
cofibrantly generated left proper Thomason model structure for
every accessible $W$. Apply D.7 and the model-category derivator
theorem cited there. Its model cofibrations preserve weak
equivalences under every cobase change, so they are intrinsic and
its first factorization gives enough intrinsic cofibrations.
The same source's Th\'eor\`eme 6.1.11 identifies properness with
generation by maps to a point; Th\'eor\`eme 5.2.15 identifies this
with model right properness. Right proper model fibrations are
intrinsic, giving the stated factorization.

For clarity the converse has a short Brown argument, which also
checks the POLES E.2 step used here. If intrinsic fibrations are
enough, their trivial members are stable under all pullbacks by
the dual proof of D.8. Every object of $\operatorname{Cat}$ is
intrinsically fibrant, since a fundamental localizer is stable
under products. Let $W^{\mathrm{asph}}$ be generated by the
$W$-aspherical categories. For a trivial intrinsic fibration
$T\to B$, its pullbacks over $B/b$ are weak equivalences with
$W$-aspherical targets, so $T/b$ is $W$-aspherical. Theorem A for
$W^{\mathrm{asph}}$ implies $T\to B\in W^{\mathrm{asph}}$.
Brown's lemma, applied to the identity functor from this category
of fibrant objects, gives $W\subset W^{\mathrm{asph}}$; the reverse
inclusion is definitional. The cited properness theorem now applies.

No initial-terminal identification is implicit: ordinary
$\operatorname{Cat}$ is not pointed. The model category of pointed
objects supplies that identification, and its finite homotopy
limits and colimits give D.6. Thus the pointed tests are correctly
typed. This result does not infer accessibility merely from a
coefficient descent axiom, nor does it impose right properness
to obtain rectification.
\end{proof}

\begin{result}{D.19}{Refuted}{Descent does not replace two out of three}
For the represented derivator of sets, the class $W^1_{\mathrm{Set}}$
defined by fully faithful restriction on locally constant objects
fails two out of three. Therefore D.7--D.9 and D.17 cannot be
applied to $(\operatorname{Cat},W^1_{\mathrm{Set}})$ merely from
its descent properties. Ordinary localization, D.2, D.4 and the
formal argument of D.15 do not require this extra hypothesis.
\end{result}
\begin{proof}
Let $i:e\to BC_2$ and $r:BC_2\to e$, so $ri=\operatorname{id}$.
Locally constant set-valued diagrams on $BC_2$ are $C_2$-sets.
The functor $r^*$ gives the trivial action and is fully faithful.
The functor $i^*$ forgets the action and is not full: on the regular
two-element $C_2$-set it sees four endofunctions, of which only two
are equivariant. Thus $r,ri\in W^1_{\mathrm{Set}}$ but
$i\notin W^1_{\mathrm{Set}}$. The multiplication and all four
functions are checked in the accompanying finite program.
This is a failure of a specific axiom, not an obstruction to forming
the relative localization in the first place.
\end{proof}

\begin{result}{D.20}{Known}{Prederivator recovery and the scope of the comparison}
The references named in Kaledin's introduction are Kevin Arlin,
\emph{A higher Whitehead theorem and the embedding of
quasicategories in prederivators}, and Daniel Fuentes-Keuthan,
Magdalena K\k edziorek and Martina Rovelli,
\emph{A model structure on prederivators for $(\infty,1)$-categories}.
Their embedding and image results are not criteria on the value
$M[W^{-1}]$ alone. They concern the family of shapes with its
transformations. In particular, homotopy categories of coherent
diagrams cannot be replaced by ordinary diagram categories in the
homotopy category.
\end{result}
\begin{proof}[Verified scope and a minimal obstruction]
\citet[Theorems 19, 29 and 37]{DArlin2020} distinguishes strict
and pseudonatural morphisms. Bicategorical full faithfulness for
small quasicategories is stated for all small category shapes;
the finite-shape assertion is conservativity for a specified
map, not the same fullness assertion. \citet[Theorems 2.16, 3.1 and 3.4 in the arXiv version]{DFKR2019} characterize
prederivators arising on the nose from quasicategories and construct
a model structure on prederivators with strict natural
transformations. This distinction from equivalence of naive
families is retained in D.15.

For the concrete obstruction take the arrows $S^0\to *$ and
$*\to S^1$. A map between these objects in the coherent arrow
category consists of the endpoint maps together with a homotopy
of the resulting square. Its mapping space is the homotopy
pullback of
\[
 *\longrightarrow \operatorname{Map}(S^0,S^1)
                  \longleftarrow S^1,
\]
where the right map is the diagonal and the left map is a fixed
constant pair. Its components form $\mathbb Z$: the long exact
sequence gives the cokernel of
$\mathbb Z\xrightarrow{\Delta}\mathbb Z^2$.
Thus there are infinitely many maps in
$\operatorname{Ho}\operatorname{Fun}([1],\mathcal S)$.
In $\operatorname{Fun}([1],\operatorname{Ho}\mathcal S)$ the
corresponding set has one element. The indexing arrow $[1]$ is
already sufficient to detect the lost homotopy of a square.
\end{proof}

\begin{result}{D.21}{Proved}{An upgrade when all Kaledin product tests are available}
Fix a category $A$. If $R_{A\times J}$ is an equivalence for every
$J\in\mathrm{Pos}$, with the stated compatible restriction
functors, then the canonical $\infty$-functor
\[
 L(M^A,W_A)\longrightarrow
       \operatorname{Fun}(NA,L(M,W))
\]
is an equivalence. The same holds using $\mathrm{Pos}^{+}$.
This is a sufficient upgrade for a specified comparison; it does
not replace $\mathrm{Pos}$ by finite posets without proof.
\end{result}
\begin{proof}
The hypothesis identifies the naive family of $(M^A,W_A)$ with
the enhanced family represented by
$X=\operatorname{Fun}(NA,L(M,W))$, since iterated coherent functor
categories satisfy the exponential identity. Apply the proof of
D.15 to this relative category. That proof shows that the
representing object $X$, with its specified map from $N(M^A)$,
is its $\infty$-localization. This is exactly the displayed
functor. Only the map with this specified provenance is covered.
\end{proof}

\begin{result}{D.22}{Open}{Exact remaining rectification and test questions}
The following targets are not settled in this package.
\begin{enumerate}
\item An intrinsic necessary-and-sufficient criterion using only
      finite-poset shapes, or only finite-category shapes, or the
      shapes of cardinality less than a prescribed regular
      $\kappa$, has not been deduced. D.15 is a criterion on
      Kaledin's $\mathrm{Pos}$ or $\mathrm{Pos}^{+}$ and cannot be
      silently restricted to any of these smaller classes.
\item The version of the prompt's mapping-space upgrade which
      assumes only the equivalences $R_{A\times J}$ for finite
      posets $J$ remains open here. D.21 establishes the larger
      shape version. The missing step is to identify the naive
      finite-shape family of $M^A$ with an enhanced family, or to
      prove the canonical $\infty$-functor is fully faithful by a
      separate finite-test theorem that applies to this source.
\item Whether Barwick--Kan fibrancy alone forces rectifiability
      at every finite poset, or at every finite category, is not
      established here. D.11 refutes the implication for all small
      shapes, already at the countable shape $B\mathbb Z$; D.12
      refutes the converse even for all small shapes.
\item Enough intrinsic cofibrations \emph{without} an initial
      object are not covered by D.8. The given least-vertex Reedy
      arguments in POLES do not themselves supply the missing
      latching coproducts for arbitrary finite diagrams. The
      initial-object version of the requested counterexample is
      impossible by D.8; the unrestricted version without that
      object is left open here.
\item The exact equivalence in IX-080 between a suitably bounded
      $\operatorname{Cat\,hot}$ representability formulation of
      (b) and its restriction-invariance formulation is not proved
      merely by a Thomason fibrant replacement. D.5 refutes the
      interpretation that only asks for a weakly equivalent
      fibrant category. The implication from the manuscript's
      path-object version of (c) to (b), in the absence of the
      comparison hypotheses in Parts B--C, remains open. D.6
      proves (c)$\Rightarrow$(a) and the actual mapping-space
      identities in a pointed $\infty$-category; D.4 proves
      rectification$\Rightarrow$(b).
\item Whether every nonaccessible fundamental localizer gives
      the unrestricted derivator on $\operatorname{Cat}$ is not
      settled. D.18 uses accessibility at the model-structure
      existence step, so its proof gives no answer after removing
      that hypothesis.
\end{enumerate}
\end{result}
\begin{proof}[What is and is not supplied by the preceding results]
For an uncountable regular $\kappa$, an ABC structure with
coproducts of fewer than $\kappa$ summands and the countable
composition axiom is a size-controlled sufficient condition for
shapes of size less than $\kappa$. The constructions in D.7 use
finite strings of arrows and countable realization stages; their
index sets remain smaller than $\kappa$. This explains the
necessary universe qualification for a model category. Merely
having finite limits and colimits corresponds to a different
assertion: a finite category with endomorphisms can have infinitely
many nondegenerate simplices in its nerve. It is not covered by
the finite-direct theorem.

D.10 rules out the proposed cubical criterion, D.11 separates
finite posets from $B\mathbb Z$, and D.13 rules out the naive use
of relative poset presentations. None is a classification of all
finite-shape rectifiable relative categories. Likewise the finite
example at IX ms.~23 fails finite completeness, not (a), (b),
(c), rectification, or enhancement. It cannot establish any of
these still-missing converses. We state these as unresolved targets
of this package, without claiming that the entire literature has
proved them independent or left them open.
\end{proof}

\begin{result}{D.23}{Proved}{The cofibration--fibration asymmetry for every homology theory}
For every nonzero spectrum $E$, the fundamental localizer
$W_E=\{u\mid E\wedge\Sigma^\infty_+Bu\text{ is an equivalence}\}$
does not have enough intrinsic fibrations. This checks the
generality of POLES C.4--C.5, including nonconnective spectra;
the sufficient cofibration and rectification statements in D.18
do not require this stronger fibration property.
\end{result}
\begin{proof}
First suppose $A\to *$ is an $E$-equivalence. The space $A$ is
nonempty and connected. Indeed, point inclusions and projection
to components split off one copy of $E$ for each component;
the fold map on two copies has nonzero kernel in some $\pi_n E$.
If $\pi_1(A,a)$ had a nonzero homomorphism to $\mathbb Z$,
divide by its image to obtain a surjection. A loop representing
$1$ and the classifying map to $S^1$ would give pointed maps
$S^1\to A\to S^1$ whose composite is homotopic to the identity.
Thus $\Sigma E$ would be a retract of the zero reduced spectrum
$E\wedge\Sigma^\infty A$, a contradiction.

Let $F=\langle x,y\rangle$ and $X=BF\simeq S^1\vee S^1$.
Every subgroup of a free group is free, as follows from its
covering graph, and every nontrivial free group surjects onto
$\mathbb Z$. Therefore every homomorphism $\pi_1(A,a)\to F$
is trivial. Since $BF$ is a $1$-type,
$\operatorname{Map}_*(A,BF)$ is contractible, and evaluation gives
$\operatorname{Map}(A,BF)\simeq BF$, with the equivalence given
by constant maps.

For any space $X$ the class
\[
 V_X=\{u\mid\operatorname{Map}(B\operatorname{cod}(u),X)
              \to\operatorname{Map}(B\operatorname{dom}(u),X)
                 \text{ is an equivalence}\}
\]
is a fundamental localizer. Two out of three and retracts are
immediate, and categories with terminal objects are contractible.
For relative Theorem A, use the natural decompositions of $BA$
and $BB$ as the homotopy colimits of their comma categories over
the common base $C$. Mapping to $X$ turns these homotopy colimits
into homotopy limits, so equivalences on every comma term give
the asserted equivalence. The preceding paragraph shows that
$V_{BF}$ contains every generator $A\to e$ of
$W_E^{\mathrm{asph}}$.

Define $\alpha:F\to F$ by
\[
 \alpha(x)=x,\qquad \alpha(y)=yxyx^{-1}y^{-1}.
\]
Its exponent-sum matrix is $I_2$, so on the reduced suspension
spectrum of the two-circle graph it is the identity matrix.
Consequently $B\alpha\in W_E$ for every spectrum $E$.
It is not surjective: the surjection
$\rho:F\to S_3$, $x\mapsto(12)$, $y\mapsto(13)$, sends both
$\alpha(x)$ and $\alpha(y)$ to $(12)$. Thus
$\rho\alpha(F)$ has two elements, whereas $\rho(F)$ has six.
The program verifies the complete $S_3$ multiplication table,
the word, and these two subgroup sizes.

If precomposition by $B\alpha$ were an equivalence of
$\operatorname{Map}(BF,BF)$, surjectivity on components would
give $r$ with $rB\alpha\simeq\operatorname{id}$. Injectivity,
applied to $(B\alpha r)B\alpha\simeq B\alpha$, would then give
$B\alpha r\simeq\operatorname{id}$. This would make $B\alpha$
a homotopy equivalence, contrary to nonsurjectivity on $\pi_1$.
Hence $B\alpha\notin V_{BF}$ and
$W_E\ne W_E^{\mathrm{asph}}$. The accessibility-free Brown
direction in the proof of D.18 shows that enough intrinsic
fibrations would force equality. This contradiction completes
the proof without using a right-properness test along an
unverified fibration.
\end{proof}
